% Shared preamble for the CSP dichotomy blueprint.
% House style follows the earlier `zhuk_centers` blueprint: numbered theorem
% environments shared across a single counter, a `longtheorem` upright variant
% for statements with many numbered hypotheses, and explicit `convention`
% environments for standing hypotheses.
\usepackage[a4paper,margin=27mm,headheight=14pt]{geometry}
\usepackage[T1]{fontenc}
\usepackage{lmodern}
\usepackage{amsmath,amssymb,amsthm}
\usepackage{longtable,array}
\usepackage{fancyhdr}
\usepackage[hidelinks]{hyperref}

\newtheoremstyle{mainplain}{9pt}{9pt}{\itshape}{}{\bfseries}{.}{0.55em}{}
\newtheoremstyle{mainroman}{9pt}{9pt}{\normalfont}{}{\bfseries}{.}{0.55em}{}
\theoremstyle{mainplain}
\newtheorem{theorem}{Theorem}[section]
\newtheorem{lemma}[theorem]{Lemma}
\newtheorem{proposition}[theorem]{Proposition}
\newtheorem{corollary}[theorem]{Corollary}
\theoremstyle{mainroman}
\newtheorem{longtheorem}[theorem]{Theorem}
\newtheorem{longlemma}[theorem]{Lemma}
\newtheorem{longproposition}[theorem]{Proposition}
\newtheorem{longcorollary}[theorem]{Corollary}
\theoremstyle{definition}
\newtheorem{definition}[theorem]{Definition}
\newtheorem{remark}[theorem]{Remark}
\newtheorem{convention}[theorem]{Convention}
\newtheorem{hazard}[theorem]{Formalization note}

\setlength{\parindent}{1.25em}\setlength{\parskip}{0pt}\setlength{\emergencystretch}{2em}
\newcommand{\alphlabels}{\renewcommand{\labelenumi}{\textup{(\alph{enumi})}}%
  \renewcommand{\theenumi}{\alph{enumi}}}
\newcommand{\romanlabels}{\renewcommand{\labelenumi}{\textup{(\roman{enumi})}}%
  \renewcommand{\theenumi}{\roman{enumi}}}
\providecommand{\tightlist}{\setlength{\itemsep}{0pt}\setlength{\parskip}{0pt}}
\numberwithin{equation}{section}
\setcounter{tocdepth}{2}\setcounter{secnumdepth}{3}
\pagestyle{fancy}\fancyhf{}
\fancyhead[L]{\small CSP dichotomy}\fancyhead[R]{\small\nouppercase{\leftmark}}
\fancyfoot[C]{\thepage}\renewcommand{\headrulewidth}{0.3pt}

% ---------------------------------------------------------------- operators
\DeclareMathOperator{\Sg}{Sg}
\DeclareMathOperator{\Clo}{Clo}
\DeclareMathOperator{\Pol}{Pol}
\DeclareMathOperator{\Inv}{Inv}
\DeclareMathOperator{\Var}{Var}
\DeclareMathOperator{\Con}{Con}
\DeclareMathOperator{\Sol}{Sol}
\DeclareMathOperator{\Expanded}{Exp}
\DeclareMathOperator{\CSP}{CSP}
\DeclareMathOperator{\GenSAT}{SAT}
\DeclareMathOperator*{\proj}{pr}
\newcommand{\ConOne}{\operatorname{Con}_1}

% ---------------------------------------------------------------- algebras
\newcommand{\bA}{\mathbf A}
\newcommand{\bB}{\mathbf B}
\newcommand{\bC}{\mathbf C}
\newcommand{\bD}{\mathbf D}
\newcommand{\bE}{\mathbf E}
\newcommand{\bR}{\mathbf R}
\newcommand{\bS}{\mathbf S}
\newcommand{\bZ}{\mathbf Z}
\newcommand{\Vn}{\mathcal V_{n}}

% ---------------------------------------------------------- subuniverse types
% The six types of Zhuk's theory.  Typed as upright small capitals so that a
% type never reads as a product of variables.
\newcommand{\TBA}{\textup{\textsc{ba}}}
\newcommand{\TC}{\textup{\textsc{c}}}
\newcommand{\TS}{\textup{\textsc{s}}}
\newcommand{\TD}{\textup{\textsc{d}}}
\newcommand{\TL}{\textup{\textsc{l}}}
\newcommand{\TPC}{\textup{\textsc{pc}}}
\newcommand{\MT}[1]{\mathcal M#1}
\newcommand{\TML}{\MT\TL}
\newcommand{\TMPC}{\MT\TPC}
\newcommand{\TMD}{\MT\TD}

% `C <_{T}^{A} B' and its relatives.
\newcommand{\subt}[2][]{<_{#2}^{#1}}
\newcommand{\subte}[2][]{\le_{#2}^{#1}}
\newcommand{\dsubt}[2][]{\dot<_{#2}^{#1}}
\newcommand{\dsubte}[2][]{\dot\le_{#2}^{#1}}
\renewcommand{\lll}{\mathrel{\lhd\!\lhd\!\lhd}}
\newcommand{\dlll}{\mathrel{\dot{\lhd}\!\dot{\lhd}\!\dot{\lhd}}}
\newcommand{\sdle}{\le_{\mathrm{sd}}}

% congruence cover  sigma^*
\newcommand{\cover}[1]{#1^{*}}
% bridge collapse
\newcommand{\bcol}[1]{\widetilde{#1}}
% absorption
\newcommand{\ab}{\lhd}
\newcommand{\abin}{\lhd_{\mathrm{bin}}}
\newcommand{\abc}{\lhd_{\mathrm{c}}}
